\chapter{Implementations and Findings}

This chapter reports the four implementations that validate the methodology of Chapter 5. It is organised by the five steps of that methodology rather than by library. Each section takes one step and reports how all four implementations met it, so that the section reads as a test of that step across four environments. Sections 6.1 to 6.5 follow the five steps in order. Section 6.6 collects the findings that recur across the cases.

The four implementations appear in the same fixed order throughout the chapter: Apache Jena in Java, rdflib.js in TypeScript, dotNetRDF in C\#, and RDFLib in Python. Table 3 sets them side by side from host library to standout limitation.

\begin{table}[h]
\centering
\caption{The four implementations from host to outcome. Repositories are at \texttt{github.com/Irfan-Ullah-cs/}.}
\label{tab:implementations-overview}
\small
\renewcommand{\arraystretch}{1.3}
\begin{tabularx}{\textwidth}{@{}>{\raggedright\arraybackslash}p{3.2cm} *{4}{>{\raggedright\arraybackslash}X}@{}}
\toprule
 & Apache Jena (Java) & rdflib.js (TypeScript) & dotNetRDF (C\#) & RDFLib (Python) \\
\midrule
Host library       & Jena 6.0.0            & rdflib.js   & dotNetRDF      & RDFLib 7.6.0 \\
UCUM backend       & systems-ucum + Indriya & ucum-lhc   & Fhir.Metrics   & ucumvert + Pint \\
SPARQL engine      & full                  & none        & full           & full \\
Operator delivery  & equality in query, rest reframed onto Java API & reframed onto store API & on the query surface & on the query surface \\
Numeric precision  & fixed-precision decimal & double precision & decimal, \textasciitilde{}28 digits & arbitrary-precision decimal \\
Tests              & 163                   & 157         & 251            & 253 \\
Repository         & \href{https://github.com/Irfan-Ullah-cs/jena-ucum}{jena-ucum} & \href{https://github.com/Irfan-Ullah-cs/rdflib.js-ucum}{rdflib.js-ucum} & \href{https://github.com/Irfan-Ullah-cs/dotnetrdf-ucum}{dotnetrdf-ucum} & \href{https://github.com/Irfan-Ullah-cs/rdflib-ucum}{rdflib-ucum} \\
\bottomrule
\end{tabularx}
\renewcommand{\arraystretch}{1}
\end{table}

Each section reports the information that determined a decision, named a mechanism, or produced a limitation. The full detail for each implementation is in the public repository.


\section{Identify Candidate Libraries}

This step looks for two things in each language: an RDF library to host the implementation and a UCUM backend to parse and compute. A short survey of each ecosystem is usually enough to see the landscape. RDF libraries can be compared on a handful of ordinary signals, none decisive alone: how long the library has existed, how widely it is adopted, how many people contribute to it, and how complete its feature set is. Completeness sits among these signals rather than above them. A library can lack a major feature and still be the natural choice, because a community has settled on it and built real systems on top of it. Adoption of that kind is its own kind of evidence. On the backend side the field in every language is narrow, and the property that matters most is whether a library can take an arbitrary UCUM expression at runtime rather than a fixed vocabulary set in advance.

For three of the four languages the survey points to one library without much contest. In Python, RDFLib has been the general-purpose RDF library of the ecosystem for over two decades, with the broadest contributor base of the hosts considered here and a full SPARQL 1.1 engine written in Python. In JavaScript, rdflib.js holds a long-established place among Linked Data applications; Section 4.2 records that it is a read-write Linked Data toolkit rather than a query engine, yet its adoption and its role in the community are what recommend it, the missing engine not withstanding. In C\#, dotNetRDF has been maintained since 2009 and ships the established in-process SPARQL 1.1 engine of the .NET world. 

Java is the exception, and deliberately so. The implementation accompanying the CDT specification was a fork of Apache Jena built against version 3, with the datatype behaviour hardcoded into the modified distribution. Jena was kept for this thesis to ask a specific question: how much of what the fork achieved from the inside can be recovered from the outside, as a standalone extension that leaves the library untouched. Carrying the work forward also meant moving from Jena 3 to the current 6.0.0. Eclipse RDF4J is the other mature Java host and nothing here counts against it; Jena was the platform the prior work sat on, and continuing that work is the point of the case.

The backends were chosen on the same mixture of capability and standing. In Java the work inherited the JSR-385 pairing already used in the fork, with systems-ucum supplying the UCUM grammar and unit definitions and Indriya the conversion and arithmetic. In JavaScript, ucum-lhc comes from the Lister Hill Center of the US National Library of Medicine and covers validation, conversion, and decomposition to SI base units on its own. In C\#, Fhir.Metrics, developed by Firely for the FHIR standard, loads the official UCUM definition file at runtime and parses expressions from the string. Python takes the same paired shape as Java. No single Python library does both: ucumvert is the only native UCUM parser, young but generated directly from the official UCUM definition file, and it computes through Pint, the established unit-arithmetic library of the ecosystem. So two of the four languages pair a UCUM-specific parser with a general arithmetic engine, ucumvert with Pint and systems-ucum with Indriya, while C\# and JavaScript each have a single library that does both.

What this step cannot settle is how far each host can be reached from outside. That is the question the next section takes up, and in Java it is the whole point of the case.


\section{Analyze Extension Points and APIs}

This step reads the internals of each host and each backend to learn what is reachable and how. The question divides three ways: how the library turns a typed literal into a value at parse time, where and whether it lets an extension reach the SPARQL operators, and what the backend offers for parsing and arithmetic. The answers vary more between the four hosts than any other step in the chapter, and that variance is the finding.

Datatype registration is the one area where every host offered a usable path, though of different kinds. RDFLib exposes a documented function, bind, that ties a datatype IRI to a value class and a constructor, after which the library calls the constructor on the lexical form at parse time and stores the result. Jena exposes a comparable public registry, TypeMapper, and calls the registered datatype's parse method when it reads a matching literal, marking the literal ill-typed if the parse throws. dotNetRDF has no single registration call. Recognition is achieved instead by subclassing the graph so that literal creation returns a quantity node for the cdt:ucum IRI, and by adding that IRI to the engine's list of recognized types so the comparison machinery treats it as a value rather than an opaque string. rdflib.js offers nothing here at all: its literal constructor performs no dispatch on the datatype IRI and calls no registered code, so a cdt:ucum literal is stored exactly like any other typed string. The consequence for rdflib.js is that recognition has to live inside the extension and be applied wherever literals enter it, the situation Chapter 5 describes for a host with no registration mechanism.

The SPARQL operators are where the hosts diverge most, and Jena is the clearest case because its engine treats its own operators inconsistently. Equality is delegated. When ARQ evaluates the equality operator on two typed literals it routes the call, through its literal layer, to the registered datatype's own equality method. Implementing that method is all native SPARQL equality requires, with no interception of any kind. Comparison and arithmetic are not delegated. The less-than and greater-than operators, the four arithmetic operators, the ORDER BY comparator, and the SUM and AVG accumulators all evaluate inside engine code that consults no registry and calls no method on the datatype; an unrecognized type raises an internal error. There is no interface to implement and no registry to populate for these, and compiled Java offers no stable way to intervene from outside. The split is therefore imposed by the engine: one obligation reachable through a clean public method, the rest reachable only by forking.

The other three hosts each sit at a different point. dotNetRDF is the most accommodating: arithmetic operators can be added through a public operator registry, equality follows from adding the datatype to the recognized-types list, the node comparer used in filters is a public settable property, and custom functions have a public factory. Only one operator path, the ORDER BY comparer, has no public setter, and it is reachable through a single reflective access to a private field, the one non-public point in the whole extension. RDFLib exposes none of these as extension points, but Python permits replacing functions at runtime, so the operators are reachable by patching, and the patching has to reach further than it first appears. The operator functions live in one module, but the parser holds its own references to them captured when it loads, and ordering and the aggregates hold further references of their own in separate modules. A change to the operator module alone would leave those paths running the original code while looking complete. rdflib.js has almost no SPARQL surface to reach: its query layer evaluates only basic patterns and equality filters, the three comparison operators all collapse to equality, and ordering, binding, and aggregates are absent. There is no function registry. What it does expose is a native query API on the store, and one method on it is the funnel the others pass through, which is the single point the extension can use.

Two structural questions from Chapter 5 resolved consistently across the hosts that have the machinery at all. Where filters and value-binding expressions exist, they share one evaluation path, so a single interception covers both. Ordering and the aggregates do not share that path; they run through separate code and must be reached separately. Jena, dotNetRDF, and RDFLib each confirm this division in their own terms, and it is the reason Chapter 5 warns against assuming one hook suffices.

The backends were read for three things: what their parser accepts and returns, whether they perform dimensional arithmetic with an error on incommensurable operands, and what numeric type they compute in. All four parse an arbitrary UCUM expression and perform dimensional arithmetic. Three of the four signal incompatible dimensions correctly, by an exception or by a failure status the extension turns into one. The fourth does not. The C\# backend treats some distinct dimensions as the same, newton and pascal among them, so an addition that should be rejected is accepted. Section 6.3 returns to this. The numeric type is where they part company, and the difference sets a ceiling no outer layer can lift. The Python backend can be configured to compute in arbitrary-precision decimal, and once a hidden fallback to floating point in the parser's transformer is corrected, it does so throughout. The C\# backend computes in the platform decimal type, roughly twenty-eight significant digits, with an overflow above that range, and like the Java backend it rounds non-terminating conversion factors to that fixed precision. The JavaScript backend computes in double-precision floating point only, as defined by IEEE 754 \cite{ieee2019_754}, with no higher-precision mode, so any precision beyond a double is lost before arithmetic begins. The Java backend keeps exact precision for integer-ratio conversions but rounds non-terminating conversion factors to a fixed decimal precision, a limitation the validation section returns to. Each backend's coverage gaps were also recorded here, and they converge on the special units of Section 3.3, which the next sections report where they surface as outcomes.

This step produced, for each host, a map from every obligation to the mechanism that can satisfy it or the absence of one. Those maps are what the feasibility assessment in the next section reasons over.

\section{Assess Feasibility}

Feasibility reads over the two maps from the previous section. One records how each host exposes the operators and the registration mechanism. The other records what each backend offers for parsing and arithmetic. The obligations of the specification guide this reading, but they do not decide it on their own. A library can be the right host even when it cannot meet every obligation, because a community has already settled on it and builds real work on top of it. That standing is part of the evidence, alongside the maps. All four candidates returned a go. What differed was what the go required.

The backend side carries one gate above the others. The backend must take an arbitrary UCUM expression at runtime and perform dimensional arithmetic on it. This gate already did work in the first step, where it removed the libraries that fix their unit vocabulary at compile time. The four chosen backends all pass it. They part company on the numeric type they compute in, and the feasibility step accepts that ceiling because no outer layer can raise it. The Python backend computes in arbitrary-precision decimal, the highest of the four. The JavaScript backend computes in double precision, the lowest.

For Jena, the go required accepting that two obligations cannot be reached from outside the library. Registration, equality, and the sameDimension function are met through public APIs. Comparison and arithmetic are not. The engine evaluates them in code that consults no registry, and compiled Java offers no way in from outside without a fork. This does not make the host unusable. Jena still parses a document carrying the datatype, still recognizes the literals, and still answers the simple queries that rest on parsing and equality.

For rdflib.js, the go rests on the same kind of reasoning at a wider scale. The library has no registration mechanism and a query layer that reduces the three comparison operators to equality, so most of the operator obligations cannot be reached inside a query at all. By an obligation count alone it would look like a poor candidate. It is not, because it holds a long-established place among Linked Data applications and is the library that community reaches for. The obligations it can meet are reframed onto its native store API, the one surface the extension can intercept, and the rest are delivered through an explicit operations layer. The backend passes the runtime gate but computes in double precision, so the go is recorded together with that ceiling.

For dotNetRDF, every obligation is reachable on the query surface, through public registries and one reflective access for the ordering comparer. One backend limitation is recorded here rather than left for later. The Fhir.Metrics dimensional check compares unit axes without their exponents. It therefore reports some distinct dimensions as the same, newton and pascal among them. The extension checks that result before adding, so an addition the specification requires to fail does not fail. This is a defect in the backend, not in the extension, and the go is taken with it documented.

For RDFLib, the go is the most complete of the four. Every obligation is met. Three are met through public APIs. The other two are met by replacing the operator functions at runtime, and those replacements are reversible and produce results the specification accepts. The backend computes in arbitrary-precision decimal once a single fallback in the parser is corrected.

All four are a go. Two reached it with every obligation on the query surface, dotNetRDF and RDFLib. Two reached it by moving obligations off that surface, Jena and rdflib.js. The rule against forking the host is what separates the two pairs. Underneath all four sits the same idea: the obligations shape the decision, but a host is chosen because it is useful and used, and a partial fit on a library a community trusts is worth more than a full fit on one nobody runs.


\section{Implementation}

This step turns the feasibility record into a working extension. The four implementations built in the same fixed order: the quantity representation first, registration second, the SPARQL hooks third, the custom function fourth, and the edge cases last. Each part depends on the one before it, because registration produces objects of the representation and the hooks evaluate those objects.

The representation is a single class that wraps the backend and exposes every operation the value space needs: parsing a lexical form and writing one back, equality, comparison, arithmetic, conversion, and the dimension test. All backend access passes through it, so the rest of the extension stays independent of the backend's API. In the C\# implementation one source file is the only place that touches Fhir.Metrics. Where the first step chose a pairing rather than a single backend, that division becomes the internal structure of the class: in Python and Java the parser produces what the arithmetic library consumes. The class realises the value space of Section 3.4, so two quantities in different units of the same dimension that denote the same value compare as equal.

Registration ties cdt:ucum and cdt:ucumunit to the representation through the mechanism the feasibility record named, after which the host parses a CDT literal into a typed object rather than an opaque string. Two matters are settled here. Registering twice must be harmless, because a host application cannot be stopped from doing it. And where the library attaches no behaviour to a datatype IRI, recognition has to live inside the extension and be applied wherever literals enter it. rdflib.js works this way, carrying its own recognition into every operation.

The SPARQL hooks follow one pattern in every implementation. A handler checks whether its operands are CDT literals. If they are, it evaluates them through the representation. If they are not, it hands the call back to the original behaviour unchanged, so the engine's treatment of every other datatype is left exactly as it was. Ill-typed literals are caught before any value operation and routed into the engine's error path, so the affected solution is dropped rather than the whole query failing. How many interception points this takes is what most separates the four. In two of them, Jena and dotNetRDF, registration alone already delivered SPARQL equality, because the engine routes equality of typed literals through the registered datatype. Python sits at the other end. The operator functions, the ordering routine, and the aggregate accumulators each held their own references, and the interception reached twenty-five points across four modules before the engine behaved consistently.

The custom function cdt:sameDimension was the simplest obligation wherever a registry existed. It extracts two quantities, asks the backend whether their dimensions match, and returns a boolean.

Implementation surfaces cases the analysis cannot show in advance, and the same five recur. Units the backend cannot parse, the special units of Section 3.3 first among them. Comparison across dimensions inside ORDER BY, which must never abort the query. Unit preservation in SUM and AVG, whose accumulators work on bare numbers and would otherwise strip the unit. Precision at the numeric extremes, where the backend's number type sets the ceiling. And the dimensionless case, where a bare number must be read as a quantity with the unit one, which more than one backend needed a dedicated branch to do. Each was handled explicitly and the chosen behaviour written down, because several return in the next step as expected failures.


\section{Validation}

This step verifies that each implementation satisfies the obligations it claims and documents what it does not. It does not let each implementation grade itself against its own tests. Before any library was tested, a single reference set of behavioral requirements was drawn up, setting out what a complete implementation should be tested for, independent of any one language. Each of the four suites is then read against that same reference, and every requirement receives a status: covered, covered but behaving differently in a documented way, not applicable to the language, or absent. Reading all four against one checklist is what lets them be compared on equal terms. The result that matters is not a count of passing tests. It is an account of conformance in which every claim is backed by a test and every limitation is visible.

The four suites differ in size, and the sizes are worth reading, though not as a score, since a larger suite is not a more conformant one. RDFLib has the most at 253 tests, no fail tests. dotNetRDF has 251, of which thirty-eight are left failing; the validation step below records which cases take which form. Jena has 163 and rdflib.js 157. What the totals hide, the distribution shows, because the shape of each suite tracks where its operators live. In RDFLib the operators run through the engine, and the query file alone holds sixty-nine tests. In rdflib.js the same obligations are reframed onto the store API, and the query file holds nine, with the weight moved to parsing and arithmetic tested on the quantity class directly. Jena has no query-level test file at all, and its largest class, sixty-seven tests, exercises the Java operations API that comparison and arithmetic were reframed onto. The reframings of Section 6.3 are visible in the counts.

Each obligation is tested at the level at which the feasibility record claims it. This distinction does real work. An obligation claimed at the query level needs a test that submits a query to the engine and reads the result, because a correct method the engine never calls passes every direct test and still fails the claim. The Python implementation tests at both levels: the quantity class on its own, and then the operators, ordering, and aggregates again through queries the engine runs. The Java implementation exercises its equality mechanism through the datatype interface but submits no query, so its query-level equality rests on Jena's documented routing rather than on a test of the extension. Making that difference visible is part of the step. Arithmetic is tested across its full type coverage, and two of the four were incomplete at the same point.

The edge cases recorded during implementation are tested next. A query over data that contains an ill-typed literal must complete with the affected solutions excluded, not abort. Ordering over mixed dimensions must complete and follow the behaviour chosen for it. The aggregates must preserve units where the implementation claims they do. A literal written out and parsed back must keep its value. Activating the extension more than once must change nothing.

What the feasibility record marked as not achievable appears here as expected failures, each with the test that fails and the reason attached, stating whether the cause is a limit inherited from a dependency or a gap in the extension. Two families recur across all four. The special units of Section 3.3, whose conversion falls outside the algebra every backend implements, and precision loss where a backend computes in a bounded numeric type. The form these take differs with the backend. Where a backend's lower precision is itself the accepted behaviour, that boundary becomes a passing test that documents the loss: rdflib.js records a value beyond double range collapsing to infinity and the sixteenth-digit case collapsing to one float. dotNetRDF records all limitations as failing tests. Its decimal backend overflows above roughly twenty-eight significant digits and rounds non-terminating conversion factors, both recorded as failing tests; the latter is why a cross-unit equality such as 3.6 km/h against 1 m/s does not hold. The Celsius offset conversion and the Fhir.Metrics dimensional defect that reports newton and pascal as one dimension are also left as failing tests with the cause attached. In all, the suite leaves thirty-eight tests failing.

The output of the step is the test suite together with the conformance account it expresses: the obligations met, the level at which each is met, the reframings, and the expected failures with their causes. That account is what lets someone who did not build the extension check its claims, and because all four are written against one reference, read the four side by side.


\section{Findings Across the Cases}

The clearest finding of the chapter is that the same obligation landed differently in each host, and for one reason: how far an extension could reach into the host from outside. That reach is set by two things. The first is the host's architecture, whether it exposes its operators through a public point or evaluates them in code an extension cannot enter. The second is the host language. A dynamic language lets an extension replace functions at runtime, which is how the Python operators were reached at twenty-five points without altering the library. A compiled language offers no such entry, so the extension is held to whatever public points exist, and where there are none, as with Jena's comparison and arithmetic, the only ways left are to fork or to reframe. The variance across the four is not noise around a single result. It is the result, and a go or no-go follows from the host's architecture and its standing rather than from counting obligations met.

The rule against forking sorts the four into three positions. RDFLib and dotNetRDF keep every obligation on the SPARQL query surface. Jena has a full SPARQL engine but its comparison and arithmetic operators are evaluated inside engine code that exposes no public extension point; those obligations are reframed onto an explicit Java API while equality remains on the query surface. rdflib.js has no complete SPARQL engine. All four stay worth implementing because they are the libraries their communities use. A partial fit on a trusted host is worth more than a complete fit on one no one runs.

Table 3 sets the four implementations side by side, from the host each ran on to the limitation each carries. All four returned a go; the table shows what the go required of each.

Where the hosts diverge, the backends converge. The gaps that recur are not host-specific but backend-specific, and the same two appear whatever the language. The special units of Section 3.3, whose conversion falls outside the algebra every backend implements, and precision loss wherever a backend computes in a bounded numeric type. The numeric ceiling is the sharpest case of the second. It is set by the backend, not the host, ranges from arbitrary-precision decimal in Python to double precision in JavaScript, and no layer above the backend can raise it.